\documentclass[11pt,a4paper]{article}
\usepackage[T1]{fontenc}
\usepackage{lmodern,amsmath,amssymb,amsthm,mathtools,booktabs,microtype}
\usepackage[margin=27mm]{geometry}
\usepackage[colorlinks=true,linkcolor=blue,citecolor=blue,urlcolor=blue]{hyperref}
\setlength{\emergencystretch}{2em}
\newtheorem{theorem}{Theorem}[section]
\newtheorem{lemma}[theorem]{Lemma}
\newtheorem{corollary}[theorem]{Corollary}
\newcommand{\Q}{\mathbb Q}
\DeclareMathOperator{\N}{N}
\title{Inverse reference functions and squared reciprocal abscissas}
\author{}\date{16 November 2026}
\begin{document}\maketitle
\begin{abstract}
Five quadratic equations in five field coordinates describe an elliptic
reference point and its translate, exposing both squared reciprocal
abscissas and three units with affine inverses. The presentation includes
the exceptional fibers of its two linear graph equations: a fifth
quadratic fixes the missing coordinate on each fiber. No nonzero guard
or auxiliary coordinate is required. The reference ratios, the two
original unit functions and the inverse chord parameter are affine
outputs. The construction works for either of the fixed translators
$R_0$ and $2R_0$ on the project curve, without requiring its abscissa
to be a square. It supplies ten rational coordinates and ten rational
quadratic equations over a quadratic field.
\end{abstract}

\section{The reference curve}
Let $k$ be a field of characteristic zero, and let $a,c\in k$ satisfy
\begin{equation}\label{eq:constants}
 \delta=a-c^2,\qquad D=a\delta,\qquad
 \Delta=(a+\delta)^2-4ac^2,\qquad ac\delta\Delta\ne0.
\end{equation}
Put
\[
 E:y^2=x^3-Dx,\qquad P=(a,ac),\qquad T_0=(0,0).
\]
The complete unit reference chart of~\cite{Unit} is
\begin{equation}\label{eq:old}
 auw=1,\qquad v\bigl(v-a(u+w)-2c\bigr)=\delta.
\end{equation}
It is isomorphic to $E\setminus\{O,T_0,P,P+T_0\}$, including
all allowed tangencies and nonzero-abscissa two-torsion points.
Its physical functions are
\begin{equation}\label{eq:physical}
\begin{aligned}
 r=\frac1{x(H)}=\frac uv,\quad
 q&=\frac1{x(H-P)}=\frac wv,\\
 S=\frac{y(H)}{x(H)}=v-c-au,\quad
 T&=\frac{y(H-P)}{x(H-P)}=aw-v+c.
\end{aligned}
\end{equation}

Define a three-coordinate inverse presentation by
\begin{equation}\label{eq:base}
 auw=1,\qquad \delta h^2+ah(u+w)+2ch=1,
 \qquad v=a(u+w)+2c+\delta h.
\end{equation}
The expression $v$ is affine; it is not an additional coordinate.

\begin{lemma}\label{lem:base}
Equations~\eqref{eq:base} are polynomially isomorphic to
\eqref{eq:old}. The maps are
\[
 h=\frac{v-a(u+w)-2c}{\delta},\qquad
 v=a(u+w)+2c+\delta h.
\]
Throughout the presentation, $u,w,h$ are units with affine inverses
\begin{equation}\label{eq:inverses}
 u^{-1}=aw,\qquad w^{-1}=au,\qquad h^{-1}=v.
\end{equation}
The reference functions satisfy
\begin{equation}\label{eq:outputs}
 r=uh,\qquad q=wh,\qquad
 S=aw+c+\delta h,\qquad T=-au-c-\delta h.
\end{equation}
\end{lemma}
\begin{proof}
The second equation of~\eqref{eq:base} says $vh=1$.
It follows that $v(v-a(u+w)-2c)=\delta vh=\delta$.
Conversely~\eqref{eq:old} gives $vh=1$ for the displayed definition
of $h$, and then the second equation of~\eqref{eq:base}.
The maps are inverse affine maps. The other assertions follow
from $auw=1$ and~\eqref{eq:physical}.
\end{proof}

\section{Five quadrics with the squared outputs}
Adjoin two coordinates $A,B$ to $(u,w,h)$ and impose, in addition
to~\eqref{eq:base}, the three quadrics
\begin{align}
 aA(v+\delta w)+2chu+h-u&=0,\label{eq:graphA}\\
 aB(v+\delta u)+2chw+h-w&=0,\label{eq:graphB}\\
 a^2\bigl(\delta^2AB-(1+4cw)A-(1+4cu)B\bigr)
       +1-2(2a+c^2)h^2&=0.\label{eq:patch}
\end{align}

\begin{theorem}\label{thm:graph}
The five quadratic equations in~\eqref{eq:base} and
\eqref{eq:graphA}--\eqref{eq:patch} describe exactly the graph
\begin{equation}\label{eq:squares}
                         A=u^2h^2=r^2,\qquad B=w^2h^2=q^2.
\end{equation}
Their projection to $(u,w,h)$ is an isomorphism. Consequently they
give a complete presentation of
$E\setminus\{O,T_0,P,P+T_0\}$ with the two squared reciprocal
abscissas as coordinates. No additional nonzero condition occurs.
\end{theorem}
\begin{proof}
Substitution of~\eqref{eq:squares} into the three equations gives
zero modulo~\eqref{eq:base}. These are universal polynomial
identities. One direct verification replaces $w$ by $(au)^{-1}$
and reduces powers of $h$ by the monic relation
\begin{equation}\label{eq:normal}
 h^2+\frac{au+2c+u^{-1}}{\delta}h-\frac1\delta=0.
\end{equation}
The certificate package gives the resulting zero remainders.

Write $L_A=v+\delta w$ and $L_B=v+\delta u$.
They cannot both vanish at a geometric point. Indeed, otherwise
$u=w$, $au^2=1$ and $v=-\delta u$. Combining $vh=1$ with
the affine expression for $v$ gives
$(a+\delta)u=-2c$. Squaring and using $au^2=1$ would give
$\Delta=0$, contrary to~\eqref{eq:constants}.

Where both coefficients are nonzero, the two linear graph equations
uniquely give~\eqref{eq:squares}. Suppose $L_A=0$.
Then $L_B\ne0$, so $B=w^2h^2$. Since $v=-\delta w$ and
$vh=1$, one has $wh=-1/\delta$ and hence $B=1/\delta^2$.
The coefficient of $A$ in~\eqref{eq:patch} is now
\[
 a^2\bigl(\delta^2B-1-4cw\bigr)=-4a^2cw\ne0.
\]
Thus the fifth quadratic uniquely recovers $A=u^2h^2$ there.
The case $L_B=0$ is symmetric: $A=1/\delta^2$ and the
remaining coefficient is $-4a^2cu\ne0$.

This argument also gives the isomorphism on coordinate rings.
The opens where $L_A$ or $L_B$ is invertible cover the base.
After recovering one square, the other linear graph coefficient
and its coefficient in~\eqref{eq:patch} have no common zero:
on the first coefficient's zero locus the latter is the displayed
unit. Thus the local inverse maps agree with the polynomial
graph~\eqref{eq:squares} and glue on the whole base.
\end{proof}

\section{Affine outputs and polynomial coordinate changes}
Every one of $u,w,h,v,S,T,A,B$ is an affine output of the new
five-coordinate presentation. Equivalently one may retain
$(S,T,h,A,B)$, using the invertible affine change
\begin{equation}\label{eq:STchange}
 u=\frac{-\delta h-T-c}{a},\qquad
 w=\frac{-\delta h+S-c}{a},\qquad
 v=S-T-\delta h.
\end{equation}
The five equations remain quadratic. In these coordinates the first
two may be written
\[
 (-\delta h-T-c)(-\delta h+S-c)=a,
 \qquad h(S-T)-\delta h^2=1.
\]

The presentation with outputs $(r,S,T,A,B)$ has polynomial recovery
from the new coordinates through $r=uh$. Conversely its inverse
chord parameter is the quadratic function
\begin{equation}\label{eq:oldinverse}
                  h=\frac{ST+c(S-T)+2\delta}{2c\delta}.
\end{equation}
Indeed the unit reference equations imply
$ST+c(T-S)+2cv+2\delta=0$ and
$h=(S-T-v)/\delta$, giving~\eqref{eq:oldinverse}.
All denominators in these coordinate changes are nonzero constants.

If $a=e^2$ for $e\in k^\times$, the additional affine function
$g=h/e$ satisfies
\[
                              g^2=rq.
\]
This statement uses no extra coordinate or guard. The reference
presentation itself does not require $a$ to be a square.

For $k=\Q(\rho)$, $\rho^2=2$, the presentation uses ten rational
coordinates and ten rational quadratic equations. The ratios $S,T$
and squared reciprocal functions $A,B$ are affine outputs, so an
existing quadratic equation involving these outputs retains its
degree after this substitution. The coefficient basis of the ten
reference equations may be changed by an invertible constant rational
matrix without changing their common zero set.

\section{Two translators on the project curve}
Let
\[
 D=11-6\rho,\qquad R_0=(1+\rho,2\rho),\qquad\rho^2=2.
\]
For the translator $P=R_0$, all the constants may be chosen integral:
\begin{equation}\label{eq:R0}
 a=1+\rho,\qquad c=4-2\rho,\qquad
 \delta=-23+17\rho,\qquad\Delta=1164-824\rho.
\end{equation}
Their norms are
\[
              \N(a)=-1,\quad\N(c)=8,\quad
              \N(\delta)=-49,\quad\N(\Delta)=-3056.
\]
In particular every nonvanishing hypothesis holds, and
$a^{-1}=\rho-1$ is integral as well.

For the translator $P=2R_0$, the constants are
\[
 a=\frac{57}{8}-\frac72\rho,\qquad
 c=-\frac{87}{41}+\frac{129}{164}\rho,\qquad
 \delta=\frac{2328-272\rho}{1681}.
\]
Here
\[
 \N(\Delta)=-\frac{8203357319351}{11574317056}\ne0,
\]
and $a=(-1+7\rho/4)^2$, so the optional function $g$ is available.
Both choices therefore satisfy the complete graph theorem.
Changing the translator in a defining source also requires its
index, height and completeness arguments to use that translator;
the algebraic theorem alone asserts no change in a defining
polynomial's variable count.

\section{Exact verification}
The package
\path{computation/elliptic_charts/inverse_reference_squares_v48/}
contains the five generic equations, their universal normal-form
checks, both quadratic-field constant specifications and an independent
finite-field check of the complete fibers. The latter exhausts the
stated small finite fields and all admissible constant pairs, testing
the exceptional graph fibers as well as the ordinary ones.
These finite checks supplement the all-field proof of
Theorem~\ref{thm:graph}; they do not replace it. No elliptic rank,
prime-pattern, height or pole-realization premise is used in this paper.

\begin{thebibliography}{9}
\bibitem{Unit}
$\mathbb Z\mathbb Q$ Diophantine Project contributors,
\emph{A three-coordinate quadratic chart for an elliptic point and its translate},
\path{papers/elliptic_charts/three_coordinate_reference_20261115_v40.tex}.
\end{thebibliography}
\end{document}
